\section{From posterior inclusion probabilities to numbers}\label{sec:methods}

Every method returns, for each region, a vector
$\pip_1,\dots,\pip_p\in[0,1]$ --- one posterior inclusion probability per
variant --- together with the truth $y_j\in\{0,1\}$ known from the simulation.
All metrics below are functions of that pair. This section defines each one and
justifies the estimator, because several of the obvious implementations are
wrong in ways that silently change conclusions.

\subsection{Average precision (the ranking metric)}\label{sec:ap}

Average precision summarises the whole precision--recall curve without
committing to a threshold, which matters here because methods' PIP scales
differ by orders of magnitude: any fixed-threshold comparison partly measures
calibration rather than ranking.

Sort the variants by $\pip$ descending. Let $k$ index the distinct PIP values
(so that all variants sharing a value form one block), $\mathrm{tp}_k$ the
cumulative number of causals at the end of block $k$, $n_k$ the cumulative
number of variants, $\mathrm{prec}_k=\mathrm{tp}_k/n_k$ and
$\mathrm{rec}_k=\mathrm{tp}_k/S$. Then
\begin{equation}
\mathrm{AP}=\sum_k\bigl(\mathrm{rec}_k-\mathrm{rec}_{k-1}\bigr)\,\mathrm{prec}_k ,
\qquad \mathrm{rec}_0=0 .
\label{eq:ap}
\end{equation}

Three implementation points are load-bearing.

\textbf{It is a step integral, not a trapezoid.} Equation~\eqref{eq:ap}
evaluates precision at the operating point reached, matching
\texttt{sklearn.metrics.average\_precision\_score}. A trapezoidal rule
interpolates between operating points that the classifier never actually
occupies and is optimistically biased.

\textbf{Ties are resolved by block, not by index.} Several methods emit many
exactly-tied PIPs. Implementing AP as ``sort, then read off precision at the
causal's index'' lets the sort's tie-break --- here genomic position --- decide
the score. Because causal variants are drawn uniformly at random this would not
introduce a constant bias but would inject pure noise into precisely the
comparisons the analysis exists to make. Collapsing tied variants into one
block gives every member the same operating point, which is the only
tie-handling that is invariant to input order.

\textbf{AP is computed per fit and then averaged, never on a merged ranking.}
Pooling all variants from all regions into one ranking and computing a single
AP is a different and incorrect quantity: it rewards a method for getting the
\emph{relative calibration across regions} right, which is not what
region-level fine-mapping is asked to do. The two disagree materially. For two
fits with $\mathrm{AP}=1$ and $\mathrm{AP}=1/3$, the correct answer is the mean
$2/3\approx0.667$, whereas the merged ranking of the same data gives $0.700$.
Both values are checked against hand computation in the test suite.

A fit with no causal variants returns \texttt{NA} rather than $0$, and NAs are
excluded from means with the count reported alongside, so a method is never
credited or penalised for a fit that posed no question.

\subsection{Thresholded discovery: FDR and its honest companion}\label{sec:fdr}

For a threshold $t$ let $\mathcal{S}_t=\{j:\pip_j\ge t\}$, $n_t=|\mathcal{S}_t|$
and $\mathrm{tp}_t=\sum_{j\in\mathcal{S}_t}y_j$. Then
\begin{equation}
\mathrm{FDR}_t=\frac{n_t-\mathrm{tp}_t}{n_t},\qquad
\text{reported jointly with } n_t .
\end{equation}
Reporting FDR alone is not meaningful: a method that never crosses $t$ has
$\mathrm{FDR}_t=0$ by convention and $n_t=0$, and would otherwise win the
comparison by abstaining. The pair $(\mathrm{FDR}_t,n_t)$ is therefore always
presented together, and $n_t$ appears in the credible-set analysis
(\S\ref{sec:cs}) as the follow-up burden.

We also compute the \textbf{honesty ratio}. A Bayesian method that emits $\pip$
values is making a falsifiable claim about its own error rate: the expected
number of false positives among $\mathcal{S}_t$ is
$\widehat{\mathrm{FP}}_t=\sum_{j\in\mathcal{S}_t}(1-\pip_j)$. Comparing that to
the realised count gives
\begin{equation}
\rho_t=\frac{\mathrm{FP}_t}{\widehat{\mathrm{FP}}_t},
\end{equation}
so $\rho_t=1$ means the method's own uncertainty statement was accurate,
$\rho_t>1$ means it was overconfident by that factor. This is stricter than
FDR: it holds the method to \emph{its own} claim rather than to an external
threshold.

\subsection{Calibration}\label{sec:calibmethod}

Two complementary summaries are used, because each is separately gameable.

\textbf{Total mass ratio.} $\sum_j \pip_j / S$. A coherent posterior over which
variants are causal should place total mass equal to the expected number of
causals. A ratio of $2$ means the method claims twice as many causal variants
as exist.

\textbf{Banded reliability.} Variants are binned by PIP into
$[0,0.1)$, $[0.1,0.5)$, $[0.5,0.9)$, $[0.9,1]$. Within a bin we accumulate the
count $n_b$, the number truly causal $c_b$, and $\sum\pip$. Calibration
compares the realised rate $c_b/n_b$ to the claimed mean $\sum\pip/n_b$; this is
the reliability diagram in Figure~\ref{fig:rel}. We report
$\mathrm{ECE}_{hi}$, the expected calibration error restricted to bins with
$\pip\ge0.1$, because the $[0,0.1)$ bin contains the overwhelming majority of
variants and an unrestricted ECE is dominated by the trivially-correct
statement that almost nothing is causal.

\textbf{Why both, and why RES.} Consider a predictor emitting $\pip_j=S/p$ for
every variant. Its mass ratio is exactly $1$ and its calibration error is
exactly $0$ --- it is perfectly calibrated and completely useless. The Murphy
decomposition of the Brier score separates this out:
\begin{equation}
\mathrm{BS}=\underbrace{\mathrm{REL}}_{\text{miscalibration}}-\underbrace{\mathrm{RES}}_{\text{resolution}}+\underbrace{\mathrm{UNC}}_{\text{base rate}},
\end{equation}
where RES measures how far the method's conditional rates depart from the base
rate. The uniform predictor has $\mathrm{REL}=0$ \emph{and} $\mathrm{RES}=0$.
Resolution is therefore the component that detects calibration achieved by
saying nothing, and it is reported in Table~\ref{tab:main-sparse-continuous} and
Figure~\ref{fig:murphy} alongside REL.

\subsection{Credible sets}\label{sec:cs}

For a target mass $\alpha$, variants are added in decreasing PIP order until the
cumulative mass first reaches $\alpha$; the resulting set size is
$|CS_\alpha|$ and coverage is the indicator that the set contains a causal
variant, averaged over fits. When a posterior's \emph{total} mass never reaches
$\alpha$ the set is the whole region and a flag \texttt{set\_reached} records
it --- this is a result, not a missing value, and suppressing it would let a
maximally diffuse posterior masquerade as a confident small set. Size and
coverage are always read together (Figure~\ref{fig:cs}): a set can be made to
cover by being large, and small by abandoning coverage.

\subsection{The replication structure and the noise floor}\label{sec:noise}

The design is a split plot. Two variance components matter:
$\sigma^2_\varepsilon$, the between-iteration variance at fixed region, and
$\sigma^2_u$, the between-region-draw variance. A factor whose levels are
compared \emph{within} the same region draws (namely $S$ and $\phi$, which are
re-simulated per scenario) sees only $\sigma^2_\varepsilon$. A factor whose
levels necessarily use \emph{different} region draws --- region size, and any
term containing it --- additionally carries $\sigma^2_u$. Those are
``whole-plot'' terms.

For a term $T$ with $\mathrm{df}_T$ degrees of freedom, the expected sum of
squares under the null is $\mathrm{df}_T\lambda_T$ with
\begin{equation}
\lambda_T=\frac{\sigma^2_\varepsilon}{20}
+\mathbf{1}\{T\ \text{whole-plot}\}\cdot\frac{m\,\sigma^2_u}{2},
\qquad m=25 ,
\end{equation}
the $20$ being the fits per cell and the $2$ the regions per size class. The
noise-corrected share is
\begin{equation}
\omega^2_T=\frac{\mathrm{SS}_T-\mathrm{df}_T\lambda_T}{\mathrm{SS}_{\text{tot}}} .
\label{eq:omega}
\end{equation}

\textbf{Why the correction is necessary.} Without it, a factor that explains
nothing still receives a share proportional to its degrees of freedom, and
whole-plot terms receive a large one. In simulation, an uncorrected null factor
recorded a $14.6\%$ share where the truth was zero; corrected, it averaged
$-0.6\%$ over replicates. Negative values are reported as computed and never
clamped: a systematically negative share is a diagnostic that
$\hat\lambda$ is too large, and hiding it would hide the diagnostic.

\textbf{Estimating $\sigma^2_u$ requires a bias correction.} The naive pooled
estimator over $J$ region draws in $K$ groups is biased low by the factor
$c=K(J-1)/(JK-1)$ --- $0.758$ at $J=4$ --- because the pooled mean absorbs part
of the between-draw signal. Dividing by $c$ restores unbiasedness; verified over
120 replicates, recovering $0.0438$ against a true $0.0400$ where the naive
estimator gave $0.0332$.

\subsection{What each ``sense'' asks}\label{sec:senses}

Three different questions are routinely conflated under ``which variables
matter''. They are kept separate throughout.

\begin{description}
\item[Sense V --- variance-explaining.] Of the variation in a metric across
cells, how much does each design factor account for, after the noise
correction of \eqref{eq:omega}? This bounds every claim below it: a noise floor
of $50\%$ means no decomposition can explain more than half.
\item[Sense D --- decision-relevant.] Does changing a factor change
\emph{which method you would choose}? Measured as the rate at which the
best-performing method flips between adjacent factor levels, split into a raw
rate and a gated rate that counts only flips whose margin survives a
cross-fitted paired test at $z=2$. The gap between them is the quantity of
interest: it is how much of the apparent structure is noise.
\item[Sense G --- diagnostic.] For a specific method, which factors hurt
\emph{it} relative to what was achievable? Measured by
$\kappa=(\text{method}-\text{baseline})/(\text{ceiling}-\text{baseline})$, the
fraction of available headroom captured, with SuSiE as baseline and
PolyFun-oracle as ceiling.
\end{description}

Sense V and Sense D can disagree entirely, and in this dataset they do. A factor
can move a metric a great deal while never changing the winner (every method
degrades together), and a factor with a small main effect can flip the winner
if it acts differently on different methods. Reporting only one of them
overstates what is known.

